\chapter{User documentation}
\label{ch:user}

This chapter provides comprehensive user documentation for the code coverage visualization feature in CodeChecker. It is designed to assist end-users in installing, operating, and effectively utilizing the feature to enhance their code quality assessment workflow.

\section{Task Description}

The code coverage visualization feature extends CodeChecker's static analysis capabilities by integrating test coverage metrics. The software enables developers to convert existing test coverage data into CodeChecker's format, store it alongside static analysis results, and visualize coverage information through an interactive web interface. This integration allows users to identify code regions with both static analysis warnings and insufficient test coverage, supporting more informed decisions about code quality and testing priorities.

\section{Target Audience}

This feature is designed for software developers, quality assurance engineers, and development teams who:
\begin{itemize}
    \item Use CodeChecker for static code analysis
    \item Maintain test suites for their codebases
    \item Need to assess the relationship between static analysis findings and test coverage
    \item Want to identify untested code regions that may contain potential defects
\end{itemize}

The feature is particularly valuable in continuous integration environments where both static analysis and test coverage metrics are regularly monitored.

\section{System Requirements}

\subsection{Minimum Requirements}

\begin{itemize}
    \item \textbf{Operating System}: Linux or macOS (the feature is compatible with CodeChecker's supported platforms)
    \item \textbf{CodeChecker}: Version compatible with the coverage feature (latest stable release recommended)
    \item \textbf{Compiler}: GCC or Clang with coverage instrumentation support (typically version 4.8+ for GCC, or any recent Clang version)
    \item \textbf{Coverage Tools}: 
    \begin{itemize}
        \item \texttt{gcov} (included with GCC/Clang)
        \item \texttt{lcov} (for processing coverage data)
    \end{itemize}
    \item \textbf{Web Browser}: Modern browser with JavaScript enabled (Chrome, Firefox, Safari, or Edge)
    \item \textbf{Database}: PostgreSQL or SQLite (as required by CodeChecker)
\end{itemize}

\subsection{Optimal Requirements}

\begin{itemize}
    \item \textbf{Operating System}: Linux (Ubuntu 20.04+ or similar) or macOS 11+
    \item \textbf{CodeChecker}: Latest version with all updates
    \item \textbf{Compiler}: GCC 9+ or Clang 10+ for improved coverage instrumentation
    \item \textbf{Memory}: 4GB RAM minimum, 8GB recommended for large codebases
    \item \textbf{Disk Space}: Sufficient space for coverage data files and database storage
    \item \textbf{Network}: Stable connection for web interface access (if using remote CodeChecker server)
\end{itemize}

\section{Installation}

The code coverage feature is integrated into CodeChecker and does not require separate installation. However, ensure that the following prerequisites are met:

\subsection{Installing Coverage Tools}

On \textbf{Ubuntu/Debian}:
\begin{lstlisting}[language=bash, caption={Installing lcov on Ubuntu/Debian}]
sudo apt-get update
sudo apt-get install lcov
\end{lstlisting}

On \textbf{macOS}:
\begin{lstlisting}[language=bash, caption={Installing lcov on macOS}]
brew install lcov
\end{lstlisting}

On \textbf{Fedora/RHEL}:
\begin{lstlisting}[language=bash, caption={Installing lcov on Fedora/RHEL}]
sudo dnf install lcov
\end{lstlisting}

\subsection{Verifying Installation}

Verify that the required tools are available:
\begin{lstlisting}[language=bash, caption={Verifying coverage tools installation}]
gcc --version
gcov --version
lcov --version
CodeChecker --version
report-converter --help
\end{lstlisting}

\section{Architectural Design Decision}
\label{sec:design-decision}

CodeChecker is fundamentally a \emph{static} analysis tool~\cite{codechecker}: it examines source code without executing it. The standard CodeChecker workflow (\texttt{log} $\rightarrow$ \texttt{analyze} $\rightarrow$ \texttt{store}) relies on a compilation database that describes how to compile each source file, but never executes the resulting binaries.

Code coverage, by contrast, is a \emph{dynamic} analysis technique. It requires compiling with instrumentation flags, executing the program (typically via a test suite), and collecting runtime data about which lines were reached. This means the compilation database used by CodeChecker's static analysis workflow is insufficient---it contains no information about how to run the software, which executables to invoke, or what test inputs to provide.

For this reason, the coverage feature does not attempt to automate the dynamic analysis steps within CodeChecker itself. Instead, the user is responsible for generating coverage data using their own build system and test infrastructure. CodeChecker's role begins at the point where coverage data already exists as an LCOV~\cite{lcov} \texttt{.info} file. The \texttt{report-converter} tool converts this file into CodeChecker's internal format, and \texttt{CodeChecker store} persists it in the database alongside static analysis results.

This separation of concerns follows the same pattern used by all other external analyzers supported by CodeChecker's \texttt{report-converter} tool (e.g., Pylint, ESLint, Cppcheck). The user runs the external tool, then converts and stores the output.

\section{Coverage Workflow Overview}

The end-to-end workflow consists of five steps, the first three of which are the user's responsibility:

\begin{figure}[H]
\centering
\begin{verbatim}
  Step 1          Step 2         Step 3          Step 4         Step 5
  gcc --coverage  ./run_tests    lcov --capture  report-conv.   CC store
  (compile)       (execute)      (collect)       -t lcov        (persist)
      |               |              |               |              |
   .gcno           .gcda        coverage.info   coverage.json   Database
  (metadata)     (runtime)      (LCOV format)   (CC format)      + Web UI
\end{verbatim}
\caption{End-to-end coverage workflow. Steps 1--3 are the user's responsibility; steps 4--5 are CodeChecker's.}
\label{fig:workflow}
\end{figure}

\begin{enumerate}
    \item \textbf{Compile with coverage flags}: Build the project with \texttt{--coverage} (GCC/Clang)
    \item \textbf{Execute tests}: Run the test suite to generate \texttt{.gcda} runtime data
    \item \textbf{Capture with lcov}: Produce a standardized \texttt{coverage.info} file
    \item \textbf{Convert}: Use \texttt{report-converter -t lcov} to produce \texttt{coverage.json}
    \item \textbf{Store and view}: Use \texttt{CodeChecker store} and browse the web UI
\end{enumerate}

\section{Step-by-Step Usage}

\subsection{Step 1: Compile with Coverage Flags}

Compile the project with the \texttt{--coverage} flag. For CMake-based projects:

\begin{lstlisting}[language=bash, caption={Building with coverage using CMake}]
mkdir build && cd build
cmake .. \
  -DCMAKE_C_FLAGS="--coverage" \
  -DCMAKE_CXX_FLAGS="--coverage" \
  -DCMAKE_EXE_LINKER_FLAGS="--coverage" \
  -DCMAKE_BUILD_TYPE=Debug \
  -DCMAKE_EXPORT_COMPILE_COMMANDS=ON
make -j$(nproc)
\end{lstlisting}

For simple single-file projects:
\begin{lstlisting}[language=bash, caption={Building a single file with coverage}]
gcc --coverage -o hello hello.c
\end{lstlisting}

The \texttt{--coverage} flag instructs the compiler to instrument the binary so it records which lines are executed at runtime. This produces \texttt{.gcno} files containing compile-time metadata.

\subsection{Step 2: Run Tests}

Execute the program or test suite to generate runtime coverage data:

\begin{lstlisting}[language=bash, caption={Running tests to generate coverage data}]
# For CMake projects with CTest
ctest --output-on-failure

# For simple programs
./hello
\end{lstlisting}

This produces \texttt{.gcda} files containing execution counts for each instrumented line.

\subsection{Step 3: Capture Coverage with LCOV}

Use \texttt{lcov} to collect the \texttt{.gcno} and \texttt{.gcda} files into a single standardized \texttt{coverage.info} file:

\begin{lstlisting}[language=bash, caption={Capturing coverage data with lcov}]
lcov --capture --directory . --output-file coverage.info \
  --ignore-errors negative
\end{lstlisting}

The \texttt{--ignore-errors negative} flag handles rare cases where multi-threaded programs produce negative execution counts.

\subsection{Step 4: Convert to CodeChecker Format}

Use the LCOV report converter to transform the \texttt{coverage.info} file into CodeChecker's \texttt{coverage.json} format:

\begin{lstlisting}[language=bash, caption={Converting LCOV data to CodeChecker format}]
report-converter -t lcov -o ./cc_reports coverage.info
\end{lstlisting}

This creates \texttt{cc\_reports/coverage/coverage.json}, which contains per-file data including covered lines, uncovered lines, and function coverage statistics extracted from the LCOV file.

\subsection{Step 5: Store and View}

Optionally run static analysis to produce reports in the same output directory, then store everything:

\begin{lstlisting}[language=bash, caption={Analyzing and storing results}]
# Optional: run static analysis
CodeChecker analyze compile_commands.json -o ./cc_reports \
  --analyzers gcc -j$(nproc)

# Start server (if not already running)
CodeChecker server --port 8001 -w ~/cc_workspace &

# Store results with coverage
CodeChecker store ./cc_reports -n my-project \
  --url http://localhost:8001/Default
\end{lstlisting}

The \texttt{CodeChecker store} command automatically detects the \texttt{coverage/coverage.json} file in the reports directory and stores the coverage data alongside any static analysis results.

\section{Output Format}

The generated \texttt{coverage.json} file uses a JSON structure mapping absolute file paths to coverage data objects:

\begin{lstlisting}[caption={Example coverage.json structure}]
{
  "/path/to/hello.c": {
    "lines": [3, 4, 6, 8, 10],
    "uncovered": [12, 13],
    "functions_found": 3,
    "functions_hit": 2
  }
}
\end{lstlisting}

Each entry contains:
\begin{itemize}
    \item \texttt{lines}: Sorted list of line numbers that were executed (from LCOV \texttt{DA:} records with execution count $> 0$)
    \item \texttt{uncovered}: Sorted list of executable line numbers that were never executed (from LCOV \texttt{DA:} records with execution count $= 0$)
    \item \texttt{functions\_found}: Total number of functions in the file (from LCOV \texttt{FNF:} record)
    \item \texttt{functions\_hit}: Number of functions that were executed at least once (from LCOV \texttt{FNH:} record)
\end{itemize}

The converter also supports the legacy format (plain list of covered line numbers) for backward compatibility with older \texttt{coverage.json} files.

\section{Visualizing Coverage in the Web Interface}

The CodeChecker web interface provides two coverage views accessible through the Statistics section.

\subsection{Coverage Statistics View}

The statistics view (\texttt{Statistics} $\rightarrow$ \texttt{Code Coverage}) displays a directory-level overview of coverage data, similar to the LCOV HTML report format:

\begin{itemize}
    \item \textbf{Directory browsing}: Navigate the project's directory structure with breadcrumb navigation
    \item \textbf{Line coverage}: Number of covered lines per directory and file, with coverage bar visualization
    \item \textbf{Function coverage}: Functions found and hit per directory and file, with coverage percentage
    \item \textbf{File search}: A search field filters the table by file or directory name for quick navigation in large projects
\end{itemize}

\begin{figure}[H]
\centering
\includegraphics[width=\textwidth,max height=0.45\textheight,keepaspectratio]{coverage_stats_top}
\caption{Coverage Statistics view showing the top-level directory listing for the Apache Xerces-C++ project. Each directory displays line coverage (with bar) and function coverage (color-coded: green $\geq$ 80\%, yellow $\geq$ 50\%, red $<$ 50\%). The header shows overall statistics: 60.3\% line coverage (40,210 of 66,649 executable lines) and 66.3\% function coverage (5,952 of 8,981 functions).}
\label{fig:stats_top}
\end{figure}

\begin{figure}[H]
\centering
\includegraphics[width=\textwidth,max height=0.45\textheight,keepaspectratio]{coverage_stats_files}
\caption{File-level view after drilling into the \texttt{datatype} directory. Individual source files are listed with their line and function coverage statistics.}
\label{fig:stats_files}
\end{figure}

\begin{figure}[H]
\centering
\includegraphics[width=\textwidth,max height=0.3\textheight,keepaspectratio]{coverage_search_filter}
\caption{The search filter in action: typing ``AbstractNumericFace'' narrows the file list to matching entries, enabling quick navigation in large projects.}
\label{fig:search_filter}
\end{figure}

\subsection{File Coverage View}

Clicking a file in the statistics view opens the file-level coverage visualization:

\begin{itemize}
    \item \textbf{Source code display}: The file's source code is shown with syntax highlighting and line numbers
    \item \textbf{Line-by-line coverage}: Each line is color-coded based on LCOV data:
    \begin{itemize}
        \item \textbf{Green background}: Line was executed during tests (covered)
        \item \textbf{Red background}: Line is executable but was never executed (uncovered)
        \item \textbf{No color}: Line is not executable (comments, blank lines, declarations)
    \end{itemize}
    \item \textbf{Coverage summary}: Header shows line coverage percentage, total executable lines, and hit lines
    \item \textbf{Searchable file selector}: An autocomplete dropdown allows searching files by name or path
\end{itemize}

The coloring is based directly on the LCOV \texttt{DA:} records stored in the database, not on heuristic guessing about which lines are executable. This ensures the visualization matches exactly what the compiler's instrumentation reported.

\begin{figure}[H]
\centering
\includegraphics[width=\textwidth,max height=0.5\textheight,keepaspectratio]{coverage_file_view}
\caption{File-level coverage view for \texttt{AbstractNumericFacetValidator.cpp}. Green lines were executed during tests (e.g., the \texttt{setPattern} and \texttt{setMaxInclusive} calls); red/orange lines are executable but uncovered (e.g., the \texttt{catch} block at line 170--173); white lines (comments, braces, blank lines) are non-executable. The header shows 49.4\% line coverage (157 of 318 executable lines) and 41.5\% function coverage. Breadcrumb navigation shows the path: top level $\rightarrow$ validators $\rightarrow$ datatype.}
\label{fig:file_view}
\end{figure}

\section{Integration with Static Analysis}

The primary value of the code coverage visualization feature lies in its integration with CodeChecker's static analysis capabilities. By viewing both static analysis reports and test coverage information in the same interface, developers can:

\begin{itemize}
    \item Identify code regions with static analysis warnings that lack adequate test coverage, prioritizing these areas for additional testing.
    \item Verify that code paths flagged by static analyzers are actually exercised by the test suite, helping to distinguish between false positives and genuine issues.
    \item Make informed decisions about code quality by considering both the presence of static analysis findings and the extent of test coverage in specific code regions.
\end{itemize}

\section{Troubleshooting}

\subsection{Common Issues and Solutions}

\begin{description}
    \item[\textbf{Issue}: Only a few files appear in coverage data]
    \textbf{Cause}: The \texttt{--coverage} flag was only applied to C files (\texttt{CMAKE\_C\_FLAGS}) but not C++ files.
    \textbf{Solution}: Set both \texttt{CMAKE\_C\_FLAGS} and \texttt{CMAKE\_CXX\_FLAGS} to \texttt{--coverage}.
    
    \item[\textbf{Issue}: Coverage data not appearing in web interface]
    \textbf{Solution}: 
    \begin{enumerate}
        \item Verify that \texttt{coverage/coverage.json} exists in the reports directory
        \item Ensure source files are also stored (run static analysis or include a report referencing the files)
        \item Check the server log for ``Could not match coverage file'' warnings
    \end{enumerate}
    
    \item[\textbf{Issue}: \texttt{lcov} reports ``Unexpected negative count'']
    \textbf{Cause}: Multi-threaded programs can produce negative counts due to race conditions in coverage counters.
    \textbf{Solution}: Add \texttt{--ignore-errors negative} to the \texttt{lcov} command, or compile with \texttt{-fprofile-update=atomic}.
    
    \item[\textbf{Issue}: No files listed in coverage visualization dropdown]
    \textbf{Solution}: 
    \begin{enumerate}
        \item Verify that coverage data was successfully stored (check server log for ``Stored coverage data for N file(s)'')
        \item Ensure you have proper permissions to view the product
    \end{enumerate}
    
    \item[\textbf{Issue}: \texttt{ERR\_INSUFFICIENT\_RESOURCES} in browser]
    \textbf{Cause}: Very large projects may exceed browser connection limits.
    \textbf{Solution}: The statistics page is optimized to use a single API call. Hard-refresh the page (\texttt{Ctrl+Shift+R}) to load the latest frontend.
\end{description}

\subsection{Security Considerations}

\begin{itemize}
    \item \textbf{File Path Exposure}: Coverage data contains absolute file paths. Ensure that sensitive path information is not exposed in shared environments.
    \item \textbf{Database Access}: Coverage data is stored in CodeChecker's database. Follow CodeChecker's security guidelines for database access and user permissions.
\end{itemize}
